Table~\ref{tab:accuracy} shows the results of the overall accuracy evaluation of \app~and the two baseline approaches.
%anomaly types
It can be seen that with \app~the top-1, top-3, and top-5 hit ratios are 0.48, 0.67, and 0.72, respectively, and the MRR is 0.58.
\app~significantly outperforms the baseline approaches in terms of all these metrics.

\begin{table}[ht]
	\newcommand{\tabincell}[2]{\begin{tabular}{@{}#1@{}}#2\end{tabular}}
	\caption{Overall Accuracy Evaluation (Best Scores are in Boldface)}
	\label{tab:accuracy}
    \small
	\centering
%	\resizebox{0.48\textwidth}{!}{
		\begin{tabular}{|c|c|c|c|c|c|}
			\hline \textbf{Metric} & \textbf{\app~} & \textbf{MonitorRank} & \textbf{Microscope}  \\
			\hline HR@1   & \textbf{0.48}  & 0.32   & 0.35  \\
			\hline HR@3   & \textbf{0.67}  & 0.40   & 0.49  \\
			\hline HR@5   & \textbf{0.72}  & 0.43   & 0.59    \\
			\hline MRR & \textbf{0.58}  & 0.37   & 0.44  \\
			\hline
		\end{tabular}
	%}
	\vspace{-3pt}
\end{table}

We also evaluate the accuracy of the three approaches for different anomaly types as shown in Table~\ref{tab:performance_algorithm_compare}.
It can be seen that \app~outperforms the two baseline approaches for all the three types of anomalies.
The advantages of \app~are the most significant for performance anomaly and the least significant for traffic anomaly.

\begin{table}[h]
	\newcommand{\tabincell}[2]{\begin{tabular}{@{}#1@{}}#2\end{tabular}}
	\caption{Accuracy Evaluation by Anomaly Types (Best Scores are in Boldface)}
	\label{tab:performance_algorithm_compare}
	\small
	\centering
%	\resizebox{0.48\textwidth}{!}{
		\begin{tabular}{|c|c|c|c|c|c|}
			\hline \textbf{Metric} & \textbf{\app~} & \textbf{MonitorRank} & \textbf{Microscope}\\
			\hline   \multicolumn{4}{|l|}{\cellcolor{gray!40} Performance Anomaly}  \\
			\hline HR@1    & \textbf{0.51} & 0.27  & 0.30   \\
			\hline HR@3    & \textbf{0.65} & 0.30  & 0.46   \\
			\hline HR@5    & \textbf{0.70} & 0.32  & 0.54   \\
			\hline MRR  & \textbf{0.61} & 0.31  & 0.39   \\
	   \hline   \multicolumn{4}{|l|}{\cellcolor{gray!40} Reliability Anomaly}  \\
			\hline HR@1    & \textbf{0.30 }& 0.26  & 0.26  \\
			\hline HR@3    & \textbf{0.56} & 0.33  & 0.44   \\
			\hline HR@5    & \textbf{0.70} & 0.35  & 0.52  \\
			\hline MRR  & \textbf{0.44} & 0.32  & 0.35  \\
			\hline   \multicolumn{4}{|l|}{\cellcolor{gray!40} Traffic Anomaly} \\
			\hline HR@1    & \textbf{0.46} & 0.32  & 0.36 \\
			\hline HR@3    & \textbf{0.55}  & 0.36  & 0.46  \\
			\hline HR@5    & \textbf{0.59} & 0.41  & 0.50  \\
			\hline MRR  & \textbf{0.54}  & 0.38  & 0.46 \\
			\hline
		\end{tabular}
	%}
	\vspace{-3pt}
\end{table}

The advantages of \app~can be explained by its specially designed mechanisms for service anomaly detection and anomaly propagation analysis.
MonitorRank considers the anomaly correlations between front-end services and back-end services and the correlations decrease with the propagation.
Therefore, it tends to choose services that are closer to the front-end service in the anomaly propagation chains.
In contrast, \app~considers the correlations between two successive service calls in the pruning.
Microscope uses the 3-sigma rule to detect anomalies and thus may produce false positives when occasional fluctuations occur.
In contrast, \app~uses more specific and sophisticated models for different types of anomalies

%The result can be explained as follows.
%For MonitorRank, its insufficiency mainly contains two factors. Firstly, MonitorRank traverses the graph with calculation of node-level metrics which are aggregated from all the interfaces of the node. So if a microservice has vast interfaces, the aggregated metrics will cover up the anomaly of the faulty interface. Secondly, as the anomaly propagates through the service call graph,  anomaly similarity between the front-end service and each node will be decreased. So traversing and sorting with the similarity will prefer to nodes which are near by the entry microservice in the anomaly propagation chains.
%For Microscope, using 3 sigma algorithm for anomaly detection will misreport the root causes for the sudden fluctuations of metrics. \app~ utilizes different models for fault localization of different anomalies, which are customized with the characteristics of each anomaly. Besides, metrics used for anomaly detection are collected from edges which represent the status of interface not microservices directly. Last but not least, \app~ takes business alarm metrics not high availability metrics to calculate the similarity to rank results, which can avoid the similarity degradation efficiently.
